	%--------------------------------------------------------------------------
%  PREÁMBULO
%--------------------------------------------------------------------------
\documentclass[12pt,a4paper,twoside]{article}
\usepackage[utf8]{inputenc}       % Si tu editor utiliza la codificación
                                    % Occidental ISO latin 1
                                    % Si utiliza Unicode pon utf8
\usepackage[german]{babel}
\usepackage{eurosym} % para el euro
\usepackage{amsmath,amssymb,dsfont,amsthm, amssymb, yhmath, tipa, eqnarray}
\usepackage{graphicx}
\usepackage{array}
\usepackage[update=true]{epstopdf}
\usepackage[a4paper,lmargin=1.5cm,rmargin=1.5cm,top=2cm,bottom=2cm]{geometry}
\usepackage[none]{hyphenat}         % Para que no corte ninguna palabra
%
\epstopdfsetup{suffix=.eps}         % Sufijo añadido a los ficheros eps
% al convertirlos a pdf

\sloppy                             % Suaviza ciertas reglas de ruptura;
% recomendable con hyphenat none

\pagestyle{plain}                   % Estilo de página por defecto;
% otra opción es empty
\setlength{\fboxsep}{3mm}           % Fija la separación entre el texto
% y el marco
\setlength{\parindent}{0pt}         % Fija el tamaño de la sangría al
% comienzo de un párrafo

\usepackage[T1]{fontenc}
\usepackage{mathpazo}
\usepackage{tabularx}
\usepackage[pdftex]{color}
\usepackage{enumerate}
\usepackage{amsmath, amsthm, amssymb, yhmath}
\usepackage{setspace}
\usepackage{multicol}
\usepackage{multirow} 
\usepackage{eurosym}
\usepackage{lastpage}
%\usepackage[usenames]{color}
\usepackage{fancyhdr} % activamos el paquete
\pagestyle{fancy} % seleccionamos un estilo
% Ponemos la línea de encabezado y pie en blanco
%\makeatletter
%\def\headrule{{\color{white}\if@fancyplain
%\let\headrulewidth\plainheadrulewidth\fi
%\hrule\@height\headrulewidth\@width\headwidth
%\vskip-\headrulewidth}}
%\makeatother
%\makeatletter
%\def\footrule{{\color{white}\if@fancyplain  %{red}
%\let\footrulewidth\plainheadrulewidth\fi
%\hrule\@height\footrulewidth\@width\headwidth
%\vskip-\footrulewidth}}
%\makeatother
%\fancyhead{} % clear all header fields
%\fancyfoot{} % clear all foot fields
%\providecommand{\abs}[1]{\lvert#1\rvert}

\fancyfoot[C]{Página \thepage \ de \pageref{LastPage}} 
%\textheight 25cm
%\spanishdecimal{.}
%--------------------------------------------------------------------------
%  CUERPO
%--------------------------------------------------------------------------
%--------------------------------------------------------------------------
%  PREÁMBULO
%--------------------------------------------------------------------------
\usepackage[utf8]{inputenc}       % Si tu editor utiliza la codificación
                                    % Occidental ISO latin 1
                                    % Si utiliza Unicode pon utf8
\usepackage[german]{babel}
\usepackage{eurosym} % para el euro
\usepackage{amsmath,amssymb,dsfont,amsthm, amssymb, yhmath, tipa, eqnarray}
\usepackage{graphicx}
\usepackage{array}
\usepackage[update=true]{epstopdf}
\usepackage[a4paper,lmargin=1.5cm,rmargin=1.5cm,top=2cm,bottom=2cm]{geometry}
\usepackage[none]{hyphenat}         % Para que no corte ninguna palabra
%
\epstopdfsetup{suffix=.eps}         % Sufijo añadido a los ficheros eps
% al convertirlos a pdf

\sloppy                             % Suaviza ciertas reglas de ruptura;
% recomendable con hyphenat none

\pagestyle{plain}                   % Estilo de página por defecto;
% otra opción es empty
\setlength{\fboxsep}{3mm}           % Fija la separación entre el texto
% y el marco
\setlength{\parindent}{0pt}         % Fija el tamaño de la sangría al
% comienzo de un párrafo

\usepackage[T1]{fontenc}
\usepackage{mathpazo}
\usepackage{tabularx}
\usepackage[pdftex]{color}
\usepackage{enumerate}
\usepackage{amsmath, amsthm, amssymb, yhmath}
\usepackage{setspace}
\usepackage{multicol}
\usepackage{multirow} 
\usepackage{eurosym}
\usepackage{lastpage}
%\usepackage[usenames]{color}
\usepackage{fancyhdr} % activamos el paquete
\pagestyle{fancy} % seleccionamos un estilo
% Ponemos la línea de encabezado y pie en blanco
%\makeatletter
%\def\headrule{{\color{white}\if@fancyplain
%\let\headrulewidth\plainheadrulewidth\fi
%\hrule\@height\headrulewidth\@width\headwidth
%\vskip-\headrulewidth}}
%\makeatother
%\makeatletter
%\def\footrule{{\color{white}\if@fancyplain  %{red}
%\let\footrulewidth\plainheadrulewidth\fi
%\hrule\@height\footrulewidth\@width\headwidth
%\vskip-\footrulewidth}}
%\makeatother
%\fancyhead{} % clear all header fields
%\fancyfoot{} % clear all foot fields
%\providecommand{\abs}[1]{\lvert#1\rvert}
\usepackage{amsmath, amssymb}
\usepackage{array}

\newcounter{ej}
\newcommand{\ejercicio}[1]{\stepcounter{ej}\par\medskip\noindent\textbf{\theej.}\ \textbf{#1}\par\smallskip}
\newcommand{\hueco}[1][2cm]{\rule{#1}{0.4pt}}

\fancyfoot[C]{Página \thepage \ de \pageref{LastPage}} 
%\textheight 25cm
%\spanishdecimal{.}
%--------------------------------------------------------------------------
%  CUERPO
%--------------------------------------------------------------------------
\begin{document}
\renewcommand*{\headrulewidth}{0pt} % Para la primera página
\renewcommand*{\footrulewidth}{0.1pt}
%\renewcommand*{\footruleskip}{3pt}

%\headrulewidth 0.4pt
%\footrulewidth 0.1 pt

%\onehalfspacing

%\title{FUNCIONES\vskip .5cm \Large{PRUEBA ESCRITA}}
%\author{Curso 4º E.S.O. Instituto Isbilya}
%\date{\normalsize {Mayo de 2012}}


%\maketitle
%\doublespacing
\begin{center}


\begin{center}
\begin{minipage}[c]{3.5cm}
\vspace{-1cm} % PARA REDUCIR EL MARGEN SUPERIOR DE LA 1ª PÁGINA

 \includegraphics [width=3cm]{logo.png} 

\vspace{-1cm} % PARA REDUCIR EL MARGEN SUPERIOR DE LA 1ª PÁGINA
\end{minipage}
\begin{minipage}[c]{13.5cm}
	\begin{small}
%\begin{center}
\begin{tabularx}{13.5cm}{lXr}
  % after \\: \hline or \cline{col1-col2} \cline{col3-col4} ...
  %\multicolumn{5}{c}{Puntuaciones}\\
     
     \multicolumn{3}{c}{\textbf{1º ESO \hspace{3 cm} MATEMÁTICAS}}\\
     % & & \\
    \cline{1-3}
 
  \cline{1-3}  \\%[0.1cm]
UNIDAD 1. DIVISIBILIDAD&  & CURSO 2025-2026\\
   \cline{1-3}  \\%0.1cm
 APELLIDOS, NOMBRE: & & 13-10-2026\\
   \cline{1-3}
 
  %\hline
  %Criterio 2.4:& & p de 1 p = & & p sobre 10\\[0.2cm]
  %\hline
\end{tabularx}
%\end{center}
\end{small}
\end{minipage}
%\vskip 3pt
%\begin{center}
%\begin{tabular*}{17cm}{l@{\extracolsep{\fill}}r}	
			%\normalsize{\textbf{I.E.S. ALARIFES RUIZ FLORINDO}}&\normalsize{1º E.S.O. MATEMÁTICAS}\\
%\hline\\
			%\normalsize{U.D. 4: FRACCIONES}&\textsc{PRUEBA ESCRITA}\\
%\hline\\
			%\large{Alumn@}:& \\
%\hline\\
%			\end{tabular*}
%	\end{center}
\end{center}

\vspace{0.5cm}
 \begin{tabularx}{17cm}{lXr}
  % after \\: \hline or \cline{col1-col2} \cline{col3-col4} ...
  %\multicolumn{5}{c}{Puntuaciones}\\
    

     \multicolumn{3}{c}{}\\
     % & & \\
  %  \cline{1-3}
 
 % \cline{1-3}  \\%[0.1cm]
\begin{tabular}{| c | c | c |}
\hline
\multicolumn{3}{ |c| }{CRITERIOS EVALUACIÓN} \\ \hline
 CE 2.1 & CE 3.1 & CE 5.2  \\ \hline 
 & &   \\
 & &  \\
 \hline
\end{tabular} &  &

 \begin{tabular}{| c |}

\multicolumn{1}{c}{ \includegraphics [width=6cm]{logo2.jpg} }


\end{tabular}
\end{tabularx}

 \end{center}

\ejercicio{[CE 2.1] Reduce estas expresiones a una sola potencia.}
\begin{align*}
3^2 \cdot 5^2 &= \hueco & 6^3 \cdot 4^3 &= \hueco \\[10pt]
10^4 : 2^4 &= \hueco & 15^6 : 3^6 &= \hueco
\end{align*}

\ejercicio{[CE 5.2] ¿Cuál es el mínimo común múltiplo de estos tres números?}
\[
\text{m.c.m.}\,(12,\ 18,\ 21) = \hueco[5cm]
\]

\ejercicio{[CE 5.2] Decide si estas afirmaciones son verdaderas (V) o falsas (F).}
\begin{itemize}\setlength{\itemsep}{6pt}
    \item El m.c.d. de dos números es el menor de sus divisores comunes. 
    \item El m.c.m. de dos números es el mayor de sus múltiplos comunes. 
    \item El m.c.d. de dos números es el mayor de sus divisores comunes. 
\end{itemize}

\ejercicio{[CE 2.1] Expresa el resultado como una sola potencia.}
\begin{align*}
7^2 \cdot 7^4 &= \hueco & 6^3 \cdot 6^2 &= \hueco \\[10pt]
5^4 : 5^3 &= \hueco & 10^6 : 10^2 &= \hueco
\end{align*}



\ejercicio{[CE 5.2] Decide si estas afirmaciones son verdaderas (V) o falsas (F).}
\begin{itemize}\setlength{\itemsep}{6pt}
    \item Un número primo siempre es divisible entre 2.
    \item Un número primo solo es divisible por él mismo y por 1. 
    \item Un número primo siempre es impar. 
\end{itemize}

\newpage

\ejercicio{[CE 3.1] ¿Qué número va en el cuadradito para que se cumpla la igualdad?}
\begin{align*}
(2^{\square})^4 &= 2^8 & \square &= \hueco[2cm] \\[10pt]
(3^{\square})^3 &= 3^{12} & \square &= \hueco[2cm] \\[10pt]
(4^{\square})^2 &= 4^6 & \square &= \hueco[2cm] \\[10pt]
(5^{\square})^5 &= 5^{25} & \square &= \hueco[2cm]
\end{align*}

\ejercicio{[CE 5.2] ¿Existe relación de divisibilidad entre estas parejas de números? Explica por qué.}

\noindent
\begin{minipage}[t]{0.55\textwidth}
\begin{itemize}\setlength{\itemsep}{10pt}
    \item 2 y 36 \quad $\rightarrow$ \hueco[5cm]
    \item 9 y 29 \quad $\rightarrow$ \hueco[5cm]
    \item 35 y 7 \quad $\rightarrow$ \hueco[5cm]
    \item 100 y 20 \quad $\rightarrow$ \hueco[5cm]
\end{itemize}
\end{minipage}%
\hfill
\begin{minipage}[t]{0.4\textwidth}
\vspace{0pt}
\footnotesize\textit{Operaciones:}
\end{minipage}
\vspace{1cm}

\ejercicio{[CE 5.2] Factoriza estos números y [CE 5.2] di cuál es su máximo común divisor.}
\begin{flalign*}
44 &= \hueco[5cm] && \\[10pt]
120 &= \hueco[5cm] && \\[10pt]
\text{m.c.d.}\,(44,\ 120) &= \hueco[5cm] &&
\end{flalign*}

\ejercicio{[CE 5.2] De la siguiente lista de números, tacha los que no sean primos.}
\[
2 \qquad 4 \qquad 7 \qquad 9 \qquad 15 \qquad 17 \qquad 21 \qquad 23
\]

\ejercicio{[CE 5.2] Marca con una X si cada número es divisible por 2, 3, 5 y 11. [CE 5.2] Después, explica el criterio de divisibilidad de cada uno.}
\vspace{1cm}

\noindent
\begin{minipage}[c]{0.45\textwidth}
\renewcommand{\arraystretch}{1.6}
\begin{tabular}{|r|*{4}{>{\centering\arraybackslash}m{0.7cm}|}}
\hline
\textbf{Número} & \textbf{2} & \textbf{3} & \textbf{5} & \textbf{11} \\
\hline
1.232 & & & & \\ \hline
11.135 & & & & \\ \hline
12.390 & & & & \\ \hline
22.222.202 & & & & \\ \hline
\end{tabular}
\end{minipage}%
\hfill
\begin{minipage}[c]{0.52\textwidth}
\begin{tabular}{@{}l l@{}}
\textbf{Criterio del 2:}  & \hueco[5cm] \\[14pt]
\textbf{Criterio del 3:}  & \hueco[5cm] \\[14pt]
\textbf{Criterio del 5:}  & \hueco[5cm] \\[14pt]
\textbf{Criterio del 11:} & \hueco[5cm] \\
\end{tabular}
\end{minipage}

\end{document}

\end{document}